\subsection{Erkennung von Dotierungsmustern bei Transistoren}
% Alt: Aufgabe 25
\Aufgabe{
   \label{Aufg25}
   Ohne das Anlegen einer Spannung ergeben sich die folgenden Verteilungen von
   p- und n-Gebieten. Um was für einen Typ von Transistor handelt es sich?

   \hspace{4ex}

   \begin{itemize}
      \begin{minipage}{.45\textwidth}
         \item [a)]
         %{\small \includesvg[width=0.9\textwidth]{Bilder/Aufgaben/FigTransistorDotierung1}}
         \includesvg[width=0.9\textwidth, pretex=\small]{Bilder/Aufgaben/FigTransistorDotierung1}
         \alttext{Das Bild zeigt vier schematische Querschnittsdarstellungen von Halbleiterbauelementen mit den Beschriftungen Source (S), Gate (G), Drain (D) und Bulk (B). Bauelement mit einem rosa-roten Bereich, der positive Ladungsträger (rote Kreise mit Pluszeichen) enthält. An den Enden befinden sich blaue Bereiche mit negativen Ladungsträgern (blaue Kreise mit Minuszeichen). Der Gate-Bereich in der Mitte ist gelb markiert.}
         \vspace{4ex}
         \item [c)]
         \includesvg[width=0.9\textwidth, pretex=\small]{Bilder/Aufgaben/FigTransistorDotierung3}\alttext{Ähnelt dem ersten Bild, aber die blauen Bereiche an Source, Gate und Drain sind größer und enthalten negative Ladungsträger. Der Bulk-Bereich ist rosa-rot mit positiven Ladungsträgern.}
      \end{minipage}
      \begin{minipage}{.45\textwidth}
         \item [b)]
         \includesvg[width=0.9\textwidth, pretex=\small]{Bilder/Aufgaben/FigTransistorDotierung2}\alttext{Ähnliches Bauelement, jedoch sind die Farben umgekehrt: Der Bulk-Bereich ist blau mit negativen Ladungsträgern, während Source und Drain rosa-rot mit positiven Ladungsträgern sind. Der Gate-Bereich bleibt gelb.}
         \vspace{4ex}
         \item [d)]
         \includesvg[width=0.9\textwidth, pretex=\small]{Bilder/Aufgaben/FigTransistorDotierung4}\alttext{Bulk-Bereich in Blau mit negativen Ladungsträgern. Source und Drain sind rosa-rot mit positiven Ladungsträgern, die sich über eine größere Fläche unter dem Gate erstrecken. Der Gate-Bereich ist wieder gelb markiert.}
      \end{minipage}
   \end{itemize}
}

\Loesung{
   \begin{itemize}
      \item[\bf a)] n-Kanal selbstsperrend (Verarmungstyp)     Siehe: Abschnitt \ref{fig:SperrbereichN-KanalMOSFET}
      \item[\bf b)] p-Kanal selbstsperrend (Verarmungstyp)     Siehe: Abschnitt \ref{fig:SperrbereichN-KanalMOSFET}
      \item[\bf c)] n-Kanal selbstleitend (Anreicherungstyp)   Siehe: Abschnitt \ref{fig:OhmscherBereichN-KanalMOSFET}
      \item[\bf d)] p-Kanal selbstleitend (Anreicherungstyp)   Siehe: Abschnitt \ref{fig:OhmscherBereichN-KanalMOSFET}
   \end{itemize}
}